% =====================================================================
%  Chương 1 — Cơ sở lý thuyết và Nguyên lý hệ thống thông minh
% =====================================================================
\chapter{Cơ sở lý thuyết và Nguyên lý Hệ thống Thông minh}
\label{ch:coso}

\section{Mã nguồn dự án và sản phẩm bàn giao}
\label{sec:masource}

Toàn bộ sản phẩm của báo cáo này là mã nguồn mở và chạy lại được từ đầu:

\begin{itemize}
  \item \textbf{Kho mã nguồn:}\\
        \url{https://github.com/Nghaiz/Assignment-01-Intelligent-System}
  \item \textbf{Ứng dụng web đã triển khai công khai:}\\
        \url{https://httm-nguyen-duy-nghia-assignment-01-intelligent-system.streamlit.app}
  \item \textbf{Đồ thị Tri thức:} nạp thật lên máy chủ Neo4j Aura
        (Kernel 5.27-aura, bản enterprise), có bản offline dự phòng.
\end{itemize}

\noindent Cấu trúc kho mã và vai trò từng thành phần được tóm tắt ở
Bảng~\ref{tab:cautruc-kho}.

\begin{table}[H]
\centering
\caption{Cấu trúc kho mã nguồn và vai trò từng thành phần}
\label{tab:cautruc-kho}
\small
\begin{tabularx}{\textwidth}{@{}L{4.3cm}X@{}}
\toprule
\textbf{Đường dẫn} & \textbf{Nội dung} \\
\midrule
\texttt{src/config.py}     & Đường dẫn, hằng số, hạt giống ngẫu nhiên, kiểu vẽ biểu đồ (nguồn duy nhất --- SSOT) \\
\texttt{src/preprocess.py} & Làm sạch dữ liệu và xây dựng biểu diễn đặc trưng \\
\texttt{src/train.py}      & Định nghĩa và huấn luyện $5+5$ mô hình \\
\texttt{src/evaluate.py}   & Bộ độ đo và bốn thực nghiệm có kiểm soát \\
\texttt{src/viz.py}        & Sinh toàn bộ 44 biểu đồ của báo cáo \\
\texttt{src/kg/ontology.py} & Bản thể học, phân tầng nguy cơ và phân khúc thị trường \\
\texttt{src/kg/catalog\_pharmacy.py} & Danh mục bán lẻ dược phẩm --- mã hàng, khoảng giá, gói chăm sóc \\
\texttt{src/kg/catalog\_urban.py}    & Hạ tầng tiện ích đô thị và phép tính tài chính vay mua nhà \\
\texttt{notebooks/}        & Hai Jupyter Notebook, đủ 22 mục theo đề bài \\
\texttt{app/}              & Ứng dụng web Streamlit --- ba trang, bốn tab mỗi trang công cụ \\
\texttt{report/}           & Mã nguồn \LaTeX{} của chính báo cáo này \\
\texttt{tests/}            & 355 bài kiểm thử tự động \\
\texttt{run\_pipeline.py}  & Một lệnh tái tạo toàn bộ kết quả, khoảng 90 giây \\
\bottomrule
\end{tabularx}
\end{table}

\section{Nguồn dữ liệu}
\label{sec:nguondulieu}

Hệ thống sử dụng hai tập dữ liệu có cấu trúc thực tế thu thập từ nền tảng
Kaggle, cả hai đều đáp ứng đủ các tiêu chí: mang ý nghĩa thực tiễn, có biến mục
tiêu rõ ràng và đủ dữ liệu để huấn luyện thuật toán học máy.

\begin{itemize}
  \item \textbf{Tập Pima Indians Diabetes} --- dữ liệu chẩn đoán y tế phục vụ
        bài toán đánh giá nguy cơ tiểu đường \emph{(Link dataset:}
        \url{https://www.kaggle.com/datasets/uciml/pima-indians-diabetes-database}\emph{)}
  \item \textbf{Tập Vietnam Housing} --- dữ liệu thuộc tính vật lý và vị trí của
        các bất động sản tại Việt Nam phục vụ bài toán định giá
        \emph{(Link dataset:}
        \url{https://www.kaggle.com/datasets/huutri148/vietnam-housing-dataset}\emph{)}
\end{itemize}

\noindent Đặc tả chi tiết của từng tập --- hiện tượng thực tế, quan sát, đặc
trưng, biến mục tiêu, loại bài toán --- được trình bày ở Mục~\ref{sec:dulieu-dia}
và Mục~\ref{sec:dulieu-hou}.

\section{Hệ thống thông minh: định nghĩa và dòng chảy thông tin}
\label{sec:dinhnghia-httm}

\subsection{Định nghĩa}

Một \textbf{hệ thống thông minh} (Intelligent System) là hệ thống nhận thông
tin từ môi trường, biểu diễn thông tin ấy dưới dạng máy tính xử lý được, học
một ánh xạ từ biểu diễn sang kết luận, rồi \emph{tác động ngược trở lại} môi
trường bằng một quyết định hoặc một khuyến nghị. Ba mệnh đề trong định nghĩa
này đều là điều kiện cần:

\begin{enumerate}
  \item \textbf{Có biểu diễn.} Không có bước chuyển từ hiện tượng sang véc-tơ
        số thì không có gì để học.
  \item \textbf{Có học từ dữ liệu.} Nếu ánh xạ được viết tay bằng luật cứng thì
        đó là hệ chuyên gia, không phải hệ thống học máy.
  \item \textbf{Có tác động trở lại.} Một mô hình chỉ in ra con số rồi dừng lại
        chưa phải là hệ thống; nó mới là một thành phần của hệ thống.
\end{enumerate}

Mệnh đề thứ ba là mệnh đề dễ bị bỏ qua nhất, và cũng chính là lý do báo cáo này
dành hẳn hai tầng cuối của Đồ thị Tri thức (Mục~\ref{sec:kg-lythuyet}) cho phần
việc ấy: mô hình trả về xác suất $0{,}545$ hoặc con số $8{,}42$ tỷ đồng, còn đồ
thị mới là thứ biến hai con số ấy thành \emph{một giỏ hàng gồm năm mặt hàng với
chi phí xác định} và \emph{một khoản trả góp $47{,}5$ triệu đồng mỗi tháng cùng
mức thu nhập tối thiểu cần có}.

\subsection{Dòng chảy thông tin sáu tầng}

Hình~\ref{fig:sodo-hethong} trình bày kiến trúc chung mà cả hai hệ thống trong
báo cáo đều tuân theo. Đây là sản phẩm bàn giao số~4 theo Mục~24 của đề bài
(Sơ đồ kiến trúc hệ thống).

\begin{figure}[H]
\centering
\includegraphics[width=0.54\textwidth]{fig_system_diagram.png}
\caption[Dòng chảy thông tin sáu tầng của một hệ thống thông minh]%
{Dòng chảy thông tin sáu tầng của một hệ thống thông minh, áp dụng đồng thời
cho hệ thống chẩn đoán tiểu đường và hệ thống định giá bất động sản}
\label{fig:sodo-hethong}
\end{figure}

\bieudo{Sơ đồ khối tuần tự (sequential block diagram). Mỗi khối là một tầng xử
lý, mũi tên một chiều thể hiện hướng đi của thông tin; không có vòng lặp ngược
vì đây là kiến trúc suy luận một lượt (single-pass inference), không phải vòng
học tăng cường.}

\nhanxet{Tầng \textsc{Representation} được đặt ở vị trí trung tâm chứ không phải
tầng \textsc{Learning Model}. Đó là một phát biểu có chủ đích của đề bài: mọi
tầng phía sau chỉ thao tác trên véc-tơ đặc trưng, nên bất kỳ thông tin nào bị
đánh mất ở tầng biểu diễn thì không tầng nào sau đó khôi phục lại được.
Mục~\ref{sec:phanchieu} là phần kiểm định thực nghiệm cho nhận định này --- và
kết quả kiểm định không đơn giản như phát biểu.}

\noindent Ánh xạ cụ thể của sáu tầng sang hai bài toán được trình bày ở
Bảng~\ref{tab:sau-tang}.

\begin{table}[H]
\centering
\caption{Ánh xạ sáu tầng kiến trúc sang hai bài toán cụ thể}
\label{tab:sau-tang}
\small
\begin{tabularx}{\textwidth}{@{}L{2.5cm}XX@{}}
\toprule
\textbf{Tầng} & \textbf{Hệ thống 1 --- Tiểu đường} & \textbf{Hệ thống 2 --- Bất động sản} \\
\midrule
Môi trường     & Bệnh nhân đến khám sàng lọc & Thị trường nhà đất Việt Nam \\
Dữ liệu vào    & 768 hồ sơ Pima Indians Diabetes & 30.229 tin đăng rao bán \\
Biểu diễn      & $\mathbf{x}\in\mathbb{R}^{12}$ (8 gốc $+$ 4 thiết kế) & $\mathbf{x}\in\mathbb{R}^{84}$ sau mã hoá one-hot \\
Mô hình học    & 5 bộ phân lớp $+$ 1 mô hình cơ sở & 5 bộ hồi quy $+$ 1 mô hình cơ sở \\
Dự đoán        & Xác suất mắc bệnh $\hat p\in[0,1]$ & Giá dự báo $\hat y$ (tỷ VNĐ) \\
Ứng dụng       & Phân tầng nguy cơ, gói chăm sóc có chi phí & Phân khúc, hạn mức vay và trả góp \\
\bottomrule
\end{tabularx}
\end{table}

\section{Biểu diễn dữ liệu --- bản chất và vai trò}
\label{sec:bieudien}

\subsection{Không gian véc-tơ đặc trưng}

Mọi thuật toán học máy truyền thống đều làm việc trên một không gian số. Cho
tập dữ liệu gồm $n$ quan sát, mỗi quan sát được viết thành một véc-tơ $d$ chiều

\begin{equation}
\mathbf{x}_i = \big[x_{i1},\, x_{i2},\, \dots,\, x_{id}\big]^{\!\top} \in \mathbb{R}^{d},
\qquad i = 1,\dots,n,
\label{eq:vector-dactrung}
\end{equation}

\noindent và toàn bộ tập dữ liệu thành ma trận đặc trưng
$\mathbf{X}\in\mathbb{R}^{n\times d}$ đi kèm véc-tơ nhãn $\mathbf{y}$. Bài toán
học có giám sát là bài toán tìm một hàm $f_{\bm\theta}:\mathbb{R}^{d}\to\mathcal{Y}$
sao cho $f_{\bm\theta}(\mathbf{x})$ xấp xỉ tốt nhãn thật.

Điểm mấu chốt: \emph{hình học của không gian $\mathbb{R}^{d}$ quyết định thuật
toán nhìn thấy gì}. Với K-Nearest Neighbors, khái niệm ``gần nhau'' là khoảng
cách Euclid

\begin{equation}
D(\mathbf{x}_a,\mathbf{x}_b)=\sqrt{\sum_{j=1}^{d}\big(x_{aj}-x_{bj}\big)^{2}}.
\label{eq:euclid}
\end{equation}

\noindent Nếu một cột có thang đo hàng trăm (Insulin, đơn vị $\mu$U/mL) đứng
cạnh một cột có thang đo phần mười (DiabetesPedigreeFunction), thì trong
công thức~\eqref{eq:euclid} cột thứ nhất lấn át hoàn toàn cột thứ hai --- không
phải vì nó quan trọng hơn về mặt y học, mà chỉ vì đơn vị đo của nó lớn hơn. Đây
là lý do bước chuẩn hoá không phải là thủ tục hình thức.

\subsection{Ba tầng chuyển hoá biểu diễn}
\label{sec:ba-tang-bieudien}

Báo cáo phân biệt rõ ba phép biến đổi thường bị gộp chung dưới tên ``tiền xử
lý'', vì Thực nghiệm~3 ở cả hai chương sau đo tác động của chúng một cách
riêng rẽ:

\begin{enumerate}
  \item \textbf{Làm sạch (cleaning).} Sửa lại những giá trị mã hoá sai --- ví
        dụ giá trị $0$ ở cột Glucose không phải nồng độ đường huyết bằng không
        mà là quy ước cho ``không đo được''. Phép này \emph{sửa lại thông tin
        sai lệch}.
  \item \textbf{Thiết kế đặc trưng (feature engineering).} Tạo cột mới từ các
        cột đã có, ví dụ $\mathrm{Glucose}\times\mathrm{BMI}$. Phép này
        \emph{đổi dạng thông tin đã có} --- nó không mang thêm dữ kiện mới nào
        vào hệ thống.
  \item \textbf{Chuẩn hoá (standardisation).} Đưa mọi cột về cùng thang đo:
        \begin{equation}
          z_{ij} = \frac{x_{ij}-\mu_j}{\sigma_j},
          \label{eq:zscore}
        \end{equation}
        với $\mu_j,\sigma_j$ là trung bình và độ lệch chuẩn của cột $j$
        \textbf{học trên riêng tập huấn luyện}. Phép này cũng chỉ
        \emph{đổi dạng}.
\end{enumerate}

Sự phân biệt giữa ``đổi dạng thông tin'' và ``thêm hoặc bớt thông tin'' chính
là chìa khoá giải thích vì sao cùng một thực nghiệm cho hai kết luận trái ngược
trên hai bài toán --- xem Mục~\ref{sec:phanchieu}.

\begin{phathien}{Rò rỉ dữ liệu qua bước tiền xử lý}
Các tham số $\mu_j$, $\sigma_j$ ở công thức~\eqref{eq:zscore}, cũng như giá trị
trung vị dùng để điền khuyết, \textbf{bắt buộc} phải được ước lượng chỉ từ tập
huấn luyện rồi áp lại y nguyên cho tập kiểm tra. Tính trên toàn bộ dữ liệu là
rò rỉ dữ liệu (data leakage)~\cite{kaufman2012}: tập kiểm tra đã gián tiếp góp phần định nghĩa
phép biến đổi mà nó sắp được đánh giá. Trong dự án này, mọi bước tiền xử lý
được gói bên trong đối tượng \texttt{Pipeline} của scikit-learn, nên
\texttt{fit} về mặt cấu trúc không thể nhìn thấy tập kiểm tra.
Mục~\ref{sec:roro} trình bày một trường hợp rò rỉ thật đã xảy ra trong dự án và
hậu quả đo được của nó.
\end{phathien}

\section{Hai lớp bài toán học có giám sát}

\subsection{Phân lớp nhị phân và hàm sigmoid}

Bài toán phân lớp nhị phân có $\mathcal{Y}=\{0,1\}$. Hồi quy Logistic mô hình
hoá \emph{xác suất} thuộc lớp dương bằng cách đưa một tổ hợp tuyến tính qua
hàm sigmoid (còn gọi là hàm logistic)~\cite{bishop2006}:

\begin{equation}
\sigma(z) = \frac{1}{1+e^{-z}}, \qquad z \in \mathbb{R},
\label{eq:sigmoid}
\end{equation}

\begin{equation}
\hat p(\mathbf{x}) \;=\; P(y=1\mid \mathbf{x})
\;=\; \sigma\!\big(\mathbf{w}^{\!\top}\mathbf{x}+b\big)
\;=\; \frac{1}{1+\exp\!\big[-(\mathbf{w}^{\!\top}\mathbf{x}+b)\big]}.
\label{eq:logistic}
\end{equation}

\noindent Hàm $\sigma$ ép mọi số thực về khoảng $(0,1)$, đơn điệu tăng, đối
xứng qua điểm $(0,\,0{,}5)$, và có đạo hàm
$\sigma'(z)=\sigma(z)\big[1-\sigma(z)\big]$ --- tính chất khiến việc tối ưu
bằng gradient trở nên gọn gàng.

Hàm mất mát tương ứng là entropy chéo nhị phân (binary cross-entropy, còn gọi
là log-loss):

\begin{equation}
\mathcal{L}_{\mathrm{log}}(\bm\theta)
= -\frac{1}{n}\sum_{i=1}^{n}
\Big[\, y_i \log \hat p_i \;+\; (1-y_i)\log\big(1-\hat p_i\big) \Big].
\label{eq:logloss}
\end{equation}

\subsection{Hồi quy và hàm mất mát bình phương sai số}

Bài toán hồi quy có $\mathcal{Y}=\mathbb{R}$. Hàm mất mát chuẩn là sai số bình
phương trung bình:

\begin{equation}
\mathcal{L}_{\mathrm{MSE}}(\bm\theta)
= \frac{1}{n}\sum_{i=1}^{n}\big(y_i - \hat y_i\big)^{2},
\qquad \hat y_i = f_{\bm\theta}(\mathbf{x}_i).
\label{eq:mse}
\end{equation}

\subsection{Hàm mất mát có điều chuẩn}

Cực tiểu hoá trần các hàm~\eqref{eq:logloss} hoặc~\eqref{eq:mse} dễ dẫn tới học
thuộc (overfitting): mô hình khớp cả nhiễu của tập huấn luyện. Điều chuẩn
(regularisation) thêm một số hạng phạt vào độ lớn tham số~\cite{hastie2009}:

\begin{equation}
\mathcal{J}(\bm\theta)
\;=\; \underbrace{\frac{1}{n}\sum_{i=1}^{n} \ell\big(y_i,\, f_{\bm\theta}(\mathbf{x}_i)\big)}_{\text{sai số trên dữ liệu}}
\;+\; \underbrace{\lambda \,\lVert \bm\theta \rVert_{2}^{2}}_{\text{số hạng điều chuẩn } L_2},
\qquad
\lVert \bm\theta \rVert_{2}^{2}=\sum_{j=1}^{d}\theta_j^{2}.
\label{eq:regularized}
\end{equation}

\noindent Tham số $\lambda>0$ điều tiết đánh đổi giữa khớp dữ liệu và giữ mô
hình đơn giản. Trong scikit-learn~\cite{scikit2011}, \texttt{LogisticRegression} mặc định dùng
phạt $L_2$ với $C=1/\lambda=1$ --- chi tiết này sẽ trực tiếp giải thích một kết
quả phản trực giác ở Mục~\ref{sec:dactrung-tot}, nơi một đặc trưng thiết kế
\emph{tốt} lại làm Hồi quy Logistic \emph{kém đi}.

\section{Nguyên lý mô hình cơ sở}
\label{sec:baseline}

Một con số độ đo đứng một mình không mang thông tin. Accuracy $0{,}79$ là tốt
hay tồi phụ thuộc hoàn toàn vào việc đoán mù đạt bao nhiêu. \textbf{Mô hình cơ
sở} (baseline) là quy tắc đơn giản nhất có thể, không học gì từ đặc trưng, đóng
vai trò mốc so sánh bắt buộc (yêu cầu R6 của đề bài):

\begin{itemize}
  \item \textbf{Phân lớp} --- luôn đoán lớp chiếm đa số
        (\texttt{DummyClassifier(strategy="most\_frequent")}).
  \item \textbf{Hồi quy} --- luôn đoán trung vị của tập huấn luyện
        (\texttt{DummyRegressor(strategy="median")}).
\end{itemize}

Ở bài toán tiểu đường, mô hình cơ sở đạt Accuracy $0{,}6547$ --- một con số
thoạt nhìn không tệ --- nhưng Recall bằng $0$, tức \emph{bỏ sót toàn bộ người
bệnh}. Đây là minh hoạ trực tiếp cho nghịch lý Accuracy trên dữ liệu lệch lớp,
và là lý do mọi kết luận trong Chương~\ref{ch:tieuduong} đều căn theo $F_1$ và
Recall chứ không căn theo Accuracy.

\section{Sáu họ thuật toán học máy sử dụng}

Năm mô hình cho mỗi bài toán được chọn theo tiêu chí \emph{khác nhau về nguyên
lý học}, để phép so sánh nói lên điều gì đó về bản chất dữ liệu chứ không chỉ
so sánh các biến thể của cùng một ý tưởng.

\begin{table}[H]
\centering
\caption{Nguyên lý học, giả định và siêu tham số của các thuật toán sử dụng}
\label{tab:thuattoan}
\scriptsize
\begin{tabularx}{\textwidth}{@{}L{2.3cm}XL{2.95cm}@{}}
\toprule
\textbf{Thuật toán} & \textbf{Nguyên lý học và giả định} & \textbf{Siêu tham số dùng} \\
\midrule
Linear Regression &
Khớp siêu phẳng cực tiểu hoá~\eqref{eq:mse}. Giả định quan hệ tuyến tính giữa đặc trưng và mục tiêu; trần hiệu năng bị chặn bởi độ mạnh tương quan tuyến tính. &
Mặc định \\
\addlinespace
Logistic Regression &
Ranh giới tuyến tính trong không gian log-odds qua~\eqref{eq:logistic}, phạt $L_2$ theo~\eqref{eq:regularized}. Nhạy với đa cộng tuyến. &
\texttt{max\_iter=1000} \\
\addlinespace
K-Nearest Neighbors &
Không học tham số~\cite{cover1967}; nhãn quyết định bởi $k$ láng giềng gần nhất theo~\eqref{eq:euclid}. \textbf{Bắt buộc} chuẩn hoá. &
\texttt{n\_neighbors=11} \\
\addlinespace
Decision Tree &
Chia đệ quy không gian đặc trưng bằng ngưỡng trên từng trục, cực tiểu hoá độ vẩn (Gini hoặc MSE). Bất biến với phép chuẩn hoá đơn điệu. &
\texttt{max\_depth=6} (phân lớp), \texttt{12} (hồi quy) \\
\addlinespace
Random Forest &
Tổng hợp nhiều cây trên mẫu bootstrap~\cite{breiman2001} cộng với chọn ngẫu nhiên tập con đặc trưng; giảm phương sai bằng trung bình hoá. &
\texttt{n\_estimators=300/200}, \texttt{max\_depth=10/18}, \texttt{min\_samples\_leaf=5} \\
\addlinespace
Support Vector Regression &
Tìm ống $\varepsilon$ mỏng nhất chứa phần lớn điểm~\cite{cortes1995}; độ phức tạp bậc hai theo số mẫu. &
\texttt{C=10.0}, \texttt{epsilon=0.2} \\
\addlinespace
XGBoost &
Cộng dồn tuần tự các cây yếu~\cite{chen2016}, mỗi cây sửa phần dư của tổ hợp trước đó; có điều chuẩn nội tại. &
\texttt{n\_estimators=300/500}, \texttt{max\_depth=4/6}, \texttt{learning\_rate=0.05} \\
\bottomrule
\end{tabularx}
\end{table}

\section{Hệ thống độ đo đánh giá}
\label{sec:dodo}

\subsection{Ma trận nhầm lẫn và bốn độ đo phân lớp}

Với bài toán nhị phân, đối chiếu nhãn thật và nhãn dự đoán cho bốn ô của
\textbf{ma trận nhầm lẫn} (Bảng~\ref{tab:confusion-schema}).

\begin{table}[H]
\centering
\caption{Ma trận nhầm lẫn --- ý nghĩa lâm sàng của từng ô}
\label{tab:confusion-schema}
\small
\begin{tabularx}{\textwidth}{@{}L{2.1cm}L{2.6cm}X@{}}
\toprule
\textbf{Ký hiệu} & \textbf{Tên} & \textbf{Ý nghĩa trong bài toán chẩn đoán} \\
\midrule
TP & Dương tính thật & Người bệnh, được phát hiện đúng \\
TN & Âm tính thật    & Người khoẻ, được kết luận đúng \\
FP & Dương tính giả  & Người khoẻ bị báo động nhầm --- tốn thêm một lần xét nghiệm \\
FN & Âm tính giả     & \textbf{Người bệnh bị bỏ sót} --- sai lầm nghiêm trọng nhất \\
\bottomrule
\end{tabularx}
\end{table}

Bốn độ đo chuẩn được định nghĩa từ bốn ô này:

\begin{align}
\text{Accuracy}  &= \frac{TP+TN}{TP+TN+FP+FN}, \label{eq:accuracy}\\[4pt]
\text{Precision} &= \frac{TP}{TP+FP}, \label{eq:precision}\\[4pt]
\text{Recall}    &= \frac{TP}{TP+FN}, \label{eq:recall}\\[4pt]
F_1 &= 2\cdot\frac{\text{Precision}\cdot\text{Recall}}
                  {\text{Precision}+\text{Recall}}. \label{eq:f1}
\end{align}

\noindent Cách đọc bốn công thức trên trong ngữ cảnh y tế:

\begin{itemize}
  \item \textbf{Accuracy}~\eqref{eq:accuracy} --- tỷ lệ đúng chung. Mất ý nghĩa
        khi hai lớp lệch nhau: ở tập Pima, đoán ``tất cả đều khoẻ'' đã đạt
        $65{,}5\%$.
  \item \textbf{Precision}~\eqref{eq:precision} --- trong những ca bị báo động,
        bao nhiêu phần trăm là báo động thật. Precision thấp nghĩa là hệ thống
        gây phiền nhiễu.
  \item \textbf{Recall}~\eqref{eq:recall} --- trong toàn bộ người bệnh, bao
        nhiêu phần trăm được tìm ra. \textbf{Đây là độ đo quan trọng nhất của
        bài toán sàng lọc}, vì FN là sai lầm đắt nhất.
  \item $F_1$~\eqref{eq:f1} --- trung bình điều hoà của hai độ đo trên. Trung
        bình điều hoà (chứ không phải trung bình cộng) phạt nặng trường hợp một
        trong hai độ đo gần $0$: mô hình cơ sở có Precision $=$ Recall $=0$ nên
        $F_1=0$, phản ánh đúng thực chất.
\end{itemize}

Ngoài bốn độ đo trên, báo cáo dùng \textbf{ROC-AUC} --- diện tích dưới đường
cong biểu diễn Recall theo tỷ lệ Dương tính giả khi ngưỡng quyết định quét từ
$1$ về $0$. ROC-AUC~\cite{fawcett2006} đo \emph{khả năng xếp hạng} của mô hình và độc lập
với ngưỡng, nên nó là độ đo duy nhất trong nhóm không thay đổi khi ta chỉnh ngưỡng
ở Thực nghiệm~1.

\subsection{Ba độ đo hồi quy}

\begin{align}
\text{MAE}  &= \frac{1}{n}\sum_{i=1}^{n}\big| y_i - \hat y_i \big|, \label{eq:mae}\\[4pt]
\text{RMSE} &= \sqrt{\frac{1}{n}\sum_{i=1}^{n}\big( y_i - \hat y_i \big)^{2}}, \label{eq:rmse}\\[4pt]
R^{2} &= 1 - \frac{\sum_{i=1}^{n}\big( y_i - \hat y_i \big)^{2}}
                  {\sum_{i=1}^{n}\big( y_i - \bar y \big)^{2}}. \label{eq:r2}
\end{align}

\noindent MAE~\eqref{eq:mae} giữ nguyên đơn vị của biến mục tiêu, nên
``MAE $=1{,}05$'' đọc thẳng được là ``sai trung bình $1{,}05$ tỷ đồng mỗi căn''.
RMSE~\eqref{eq:rmse} bình phương sai số trước khi lấy trung bình nên phạt nặng
các ca sai lớn; khoảng cách giữa RMSE và MAE do đó là một chỉ báo về mức độ tồn
tại của các ca sai cực đoan. $R^{2}$~\eqref{eq:r2} là tỷ lệ phương sai của biến
mục tiêu được mô hình giải thích: $R^{2}=1$ là dự báo hoàn hảo, $R^{2}=0$ là
ngang với việc luôn đoán giá trị trung bình, và $R^{2}<0$ nghĩa là mô hình còn
tệ hơn cả đoán trung bình.

\begin{phathien}{$R^2$ không so sánh được giữa hai tập dữ liệu khác nhau}
Mẫu số của~\eqref{eq:r2} là phương sai của chính biến mục tiêu. Nếu một tập dữ
liệu có mục tiêu bị thu hẹp dải giá trị --- đúng trường hợp của tập bất động
sản, xem Mục~\ref{sec:price-cat-bien} --- thì mẫu số nhỏ đi và cùng một sai số
tuyệt đối sẽ cho $R^{2}$ thấp hơn. Do đó $R^{2}=0{,}61$ ở đây khắt khe hơn con
số $0{,}61$ trên một tập có dải giá trải rộng.
\end{phathien}

\section{Đồ thị Tri thức}
\label{sec:kg-lythuyet}

\subsection{Giới hạn cố hữu của mô hình hộp đen}

Một mô hình học máy được huấn luyện tốt trả lời được câu hỏi ``bao nhiêu'' và
``khả năng nào'', nhưng không trả lời được ``vì vậy thì phải làm gì''. Xác suất
$0{,}545$ không tự nó nói cho bác sĩ biết cần chỉ định xét nghiệm nào, cũng
không nói cho người bệnh biết phải mua máy đo nào và tốn bao nhiêu; con số
$8{,}42$ tỷ không tự nó nói cho người mua biết ngân hàng cho vay tối đa bao
nhiêu và mỗi tháng phải trả bao nhiêu. Khoảng trống giữa \emph{con số dự báo} và
\emph{hành động cụ thể} là khoảng trống mà tri thức miền phải lấp, và Đồ thị Tri
thức là cấu trúc dữ liệu để mã hoá tri thức ấy dưới dạng máy suy luận được.

\subsection{Định nghĩa hình thức}

Một \textbf{Đồ thị Tri thức}~\cite{hogan2021,ji2022} là bộ ba

\begin{equation}
\mathcal{KG} \;=\; \big(V,\; E,\; R\big),
\label{eq:kg}
\end{equation}

\noindent trong đó $V$ là tập \emph{thực thể} (đỉnh), $R$ là tập \emph{loại
quan hệ}, và $E\subseteq V\times R\times V$ là tập \emph{cạnh có nhãn}. Mỗi
phần tử của $E$ là một \textbf{bộ ba tri thức} (knowledge triplet)

\begin{equation}
\big(h,\; r,\; t\big) \in E,
\qquad h,t \in V,\; r \in R,
\label{eq:triplet}
\end{equation}

\noindent đọc là ``chủ thể $h$ có quan hệ $r$ với đối tượng $t$''. Ví dụ thật
lấy từ đồ thị của hệ thống thứ nhất:

\begin{center}
\small
\begin{tabular}{@{}l@{}}
\texttt{(Bệnh nhân \#192, HAS\_BIOMARKER, Glucose = 159 mg/dL)}\\
\texttt{(Mô hình K-Nearest Neighbors, PRODUCES, Xác suất rủi ro 54.5\%)}\\
\texttt{(Xác suất rủi ro 54.5\%, MAPS\_TO\_DISEASE, ICD-10 E11)}\\
\texttt{(Nhóm Nguy cơ Cao, RECOMMENDS\_PACKAGE, Gói Theo dõi Chuyên sâu)}\\
\texttt{(Gói Theo dõi Chuyên sâu, FULFILLED\_BY, Nhà thuốc FPT Long Châu)}\\
\end{tabular}
\end{center}

Năm bộ ba trên nối liền một mạch từ số đo sinh học tới nơi mua được hàng ---
năng lực mà mô hình đơn lẻ không có. Đồ thị của dự án này có $|V|=117$,
$|E|=133$ và $|R|=28$ trên 6 ca mẫu, đã nạp thật lên máy chủ Neo4j Aura.

\subsection{Kiến trúc năm tầng}
\label{sec:kg-5-tang}

Bộ ba $(V,E,R)$ ở~\eqref{eq:kg} mô tả \emph{cấu trúc}, nhưng không nói gì về
\emph{tổ chức}. Đồ thị của dự án này được tổ chức thành năm tầng ở cả hai miền
--- Hình~\ref{fig:kg-arch-medical} và Hình~\ref{fig:kg-arch-property}. Mỗi tầng
là một nhóm loại thực thể, và thông tin chỉ chảy theo một chiều từ tầng trên
xuống tầng dưới.

\begin{figure}[H]
\centering
\includegraphics[width=\textwidth]{fig_kg_architecture_medical.png}
\caption[Kiến trúc năm tầng của Đồ thị Tri thức Y sinh]%
{Kiến trúc năm tầng của Đồ thị Tri thức Y sinh --- từ chỉ số sinh tồn tới gói
chăm sóc mua được tại nhà thuốc}
\label{fig:kg-arch-medical}
\end{figure}

\begin{figure}[H]
\centering
\includegraphics[width=\textwidth]{fig_kg_architecture_property.png}
\caption[Kiến trúc năm tầng của Đồ thị Tri thức Đô thị -- Tài chính]%
{Kiến trúc năm tầng của Đồ thị Tri thức Đô thị -- Tài chính --- từ thuộc tính
căn nhà tới hạn mức vay và tiện ích hạ tầng quanh đó}
\label{fig:kg-arch-property}
\end{figure}

\bieudo{Sơ đồ tầng (layered schema diagram). Khác với ảnh chụp đồ thị của một
ca cụ thể ở Hình~\ref{fig:kg-medical} và Hình~\ref{fig:kg-property} --- vốn cho
thấy \emph{dữ liệu} --- hai hình trên cho thấy \emph{lược đồ}: đồ thị có mấy
tầng, mỗi tầng chứa loại thực thể nào, thông tin chảy theo hướng nào.}

\nhanxet{Hai kiến trúc song song nhau đến từng tầng, và đó là thiết kế có chủ
đích. Tầng 1--2 là \emph{quan sát}, tầng 3 là \emph{suy luận của mô hình}, tầng
4 là \emph{gói giải pháp có chi phí}, tầng 5 là \emph{hạ tầng cung ứng}. Nhờ
cùng khung, một câu hỏi đặt ra ở miền này luôn có câu hỏi đối ứng ở miền kia:
``gói chăm sóc tốn bao nhiêu mỗi tháng'' đối ứng với ``khoản vay trả góp bao
nhiêu mỗi tháng''.

Hai hình trên được vẽ tự động từ chính khai báo \texttt{KG\_LAYERS} trong
\texttt{src/kg/ontology.py}, nên chúng không thể mô tả một kiến trúc mà mã nguồn
không còn dùng. Một bài kiểm thử canh đúng điều đó: mọi loại thực thể mà mã
nguồn thật sự sinh ra đều phải thuộc đúng một tầng.}

\subsection{Ba vai trò của đồ thị trong hệ thống}

\begin{enumerate}
  \item \textbf{Diễn giải (interpretability).} Đường đi từ đỉnh ``Bệnh nhân''
        tới đỉnh ``Gói chăm sóc'' là một chuỗi bước đọc được bằng ngôn ngữ tự
        nhiên, khác hẳn trọng số của một cây tăng cường.
  \item \textbf{Nối tri thức miền.} Bảng phân loại ICD-10~\cite{icd10}, ngưỡng BMI của WHO~\cite{who2000},
        danh mục hàng của nhà thuốc, tỷ lệ LTV và ngưỡng DTI của ngân hàng ---
        những tri thức này không có trong dữ liệu huấn luyện và cũng không nên
        học từ dữ liệu; chúng được khai báo tường minh thành đỉnh và cạnh.
  \item \textbf{Sinh hành động có chi phí.} Truy vấn đồ thị theo tầng nguy cơ
        hoặc phân khúc giá trả về ngay một gói kèm con số: chi phí ban đầu, chi
        phí duy trì hằng tháng, hoặc khoản trả góp và mức thu nhập tối thiểu.
        Đây chính là đầu ra của tầng \textsc{Application} ở
        Hình~\ref{fig:sodo-hethong}.
\end{enumerate}

\subsection{Ranh giới trung thực về nguồn số liệu}
\label{sec:ranh-gioi-so-lieu}

Hai tầng cuối của mỗi đồ thị chứa số liệu \emph{không} đến từ tập dữ liệu huấn
luyện, nên nguồn gốc của chúng phải được nói rõ chứ không được trình bày lẫn
với kết quả đo được:

\begin{itemize}
  \item \textbf{Tên mặt hàng là tên thật} đang được bán tại hệ thống nhà thuốc.
        \textbf{Giá là khoảng tham khảo có ghi ngày chốt} (26/08/2026), biên
        soạn tay từ danh mục công khai --- không phải bản ghi giá trực tiếp. Mỗi
        mặt hàng mang theo một đường dẫn tra cứu để người dùng đối chiếu giá
        hiện hành.
  \item \textbf{Cụm tiện ích đô thị được mô hình hoá theo phân khúc}, không
        phải kết quả khảo sát thực địa từng quận: tập dữ liệu gốc chỉ có địa chỉ
        hành chính, không có toạ độ.
  \item \textbf{Điều kiện vay là minh hoạ} theo mặt bằng lãi suất công bố, không
        phải chào giá của một ngân hàng cụ thể. Ngược lại, \emph{phép tính} trả
        góp là toán học thuần tuý và kiểm chứng được --- xem
        Mục~\ref{sec:cong-thuc-nien-kim}.
\end{itemize}

Phương án thu thập giá trực tiếp từ website bán lẻ đã được cân nhắc và bị loại,
vì nó phá vỡ yêu cầu R14 về tính tái lập: giá đổi theo ngày nên hai lần chạy
cách nhau một tuần sẽ cho hai bộ số khác nhau, trong khi cả báo cáo được xây
trên cam kết ``chạy lại cho kết quả trùng khít''.

\subsection{Phép tính tài chính}
\label{sec:cong-thuc-nien-kim}

Khoản trả góp cố định hằng tháng của một khoản vay niên kim là

\begin{equation}
M \;=\; P \cdot \frac{r}{1-(1+r)^{-n}},
\qquad r=\frac{r_{\text{năm}}}{12},\quad n = 12\,T,
\label{eq:nienkim}
\end{equation}

\noindent với $P$ là dư nợ gốc, $r_{\text{năm}}$ là lãi suất danh nghĩa theo
năm và $T$ là kỳ hạn tính bằng năm. Tổng tiền lãi phải trả trong toàn bộ kỳ hạn
là $Mn - P$. Mức thu nhập tối thiểu suy ra từ ngưỡng tỷ lệ nợ trên thu nhập
(Debt-To-Income):

\begin{equation}
I_{\min} \;=\; \frac{M}{\text{DTI}_{\max}}.
\label{eq:dti}
\end{equation}

Ba công thức này được kiểm thử tự động bằng một phép tính độc lập viết lại tay,
chứ không đối chiếu với chính hàm cần kiểm.

\section{Giao thức thực nghiệm chung}
\label{sec:giaothuc}

Hai chương sau dùng chung một giao thức, để mọi so sánh giữa hai bài toán là so
sánh công bằng.

\subsection{Chia dữ liệu ba phần}

Dữ liệu được cắt làm ba, theo đúng thứ tự sau:

\begin{enumerate}
  \item \textbf{Tập ngoại kiểm (OOD), $10\%$} --- cắt ra \emph{trước tiên} và
        không tham gia bất kỳ bước nào của quá trình phát triển, kể cả tinh
        chỉnh siêu tham số.
  \item \textbf{Tập kiểm tra, $20\%$ phần còn lại} --- dùng để so sánh mô hình.
  \item \textbf{Tập huấn luyện, phần còn lại} --- dùng để học tham số.
\end{enumerate}

Lý do tách riêng tập ngoại kiểm: sau nhiều vòng thử nghiệm, tập kiểm tra thông
thường đã \emph{gián tiếp} ảnh hưởng tới lựa chọn mô hình, nên nó không còn là
ước lượng không chệch của hiệu năng thực tế. Tập ngoại kiểm đóng vai trò ``dữ
liệu thế giới thực'' chưa từng bị nhìn thấy.

\subsection{Bốn thực nghiệm có kiểm soát}

Mỗi bài toán chạy bốn thực nghiệm cùng khuôn (đề bài yêu cầu tối thiểu ba ---
R8), mỗi thực nghiệm chỉ đổi \emph{một} biến và giữ nguyên mọi thứ còn lại:

\begin{table}[H]
\centering
\caption{Bốn thực nghiệm có kiểm soát, áp dụng song song cho hai bài toán}
\label{tab:bon-thucnghiem}
\small
\begin{tabularx}{\textwidth}{@{}L{0.8cm}L{3.5cm}XL{2.3cm}@{}}
\toprule
\textbf{TN} & \textbf{Biến được thay đổi} & \textbf{Câu hỏi trả lời} & \textbf{Kết quả ở} \\
\midrule
1 & Đối đầu 5 thuật toán & Nguyên lý học nào hợp với dữ liệu này? & \ref{sec:tn1-dia}, \ref{sec:tn1-hou} \\
2 & Siêu tham số cốt lõi & Cực trị nằm ở đâu, có chạm biên dải khảo sát không? & \ref{sec:tn2-dia}, \ref{sec:tn2-hou} \\
3 & Cách biểu diễn dữ liệu & Biểu diễn hay thuật toán chi phối kết quả? & \ref{sec:tn3-dia}, \ref{sec:tn3-hou} \\
4 & Tập đánh giá (kiểm tra $\to$ ngoại kiểm) & Mô hình tổng quát hoá được tới đâu? & \ref{sec:tn4-dia}, \ref{sec:tn4-hou} \\
\bottomrule
\end{tabularx}
\end{table}

\subsection{Tính tái lập}
\label{sec:tailap}

Hạt giống ngẫu nhiên cố định tại \texttt{RANDOM\_STATE = 42} khai báo một lần
duy nhất trong \texttt{src/config.py} và được mọi thành phần dùng lại. Một lệnh
\texttt{python run\_pipeline.py} sinh lại toàn bộ mô hình, biểu đồ, bảng số và
đồ thị tri thức trong khoảng 90 giây.

Một lưu ý trung thực về mức độ tái lập: các tệp \texttt{models/*.joblib} chạy
lại cho kết quả trùng khít, nhưng vài tệp CSV của bài toán bất động sản có thể
lệch ở chữ số cuối của số thực (bậc $10^{-16}$) --- ví dụ MAE của Random Forest
đổi từ $1{,}0818214861905244$ thành $1{,}0818214861905247$. Nguyên nhân là phép
cộng dồn song song trong rừng ngẫu nhiên không có tính kết hợp. Ở độ chính xác
báo cáo --- bốn chữ số thập phân --- mọi con số đều trùng khớp.
